\section{Appendix \label{appendix}}
\subsection{Interview with Participating Musician 1 (Tomohito Yoshijima, December 15, 2017)}

\begin{description}
 \item[Impressions after watching the video and listening to the audio of this work]\mbox{}\\

 There is no precedent for music with such scores. It was difficult to create, so I didn't know what it would be like until it was finished. When it became a video, it surprisingly turned into music. Then I felt like I could get more ideas from watching the video.\\

 --- {\bf What was impressive during the recording was that you said, ``I don't understand what's happening while I'm playing.''}\\

 The understanding from the first playback to the second playback was fast (because I was listening to the playback). We were able to discuss in the booth about ``let's do it this way.'' Once it was completed and the form was decided, there was a feeling that we could control something, like a photo of a completed plastic model, or ``we're going to make this now'' - that kind of understanding emerged. The score was like an instruction manual, or rather, in recording, the playback of what we performed became like an instruction manual.\\

 --- {\bf Was it at the level of not understanding the content of the booklet, or was it the category of understanding the booklet as a booklet but not knowing until you tried it?}\\

 I could simply read the score. But I didn't know how it would sound when three people came together, or what kind of groove it would become. I think groove is the most difficult thing now, because you can't write groove in a score, like fluctuations. That's where my attention goes.\\

--- {\bf Also, drummers are generally not thought of as harmonic or melodic instruments in the world, but on SNS, Yoshijima-san occasionally expresses concern to composers about what the groove is like, and how to react to the surrounding sounds. You want them to care not just about ``what's happening?'' and improvisation, but also about rhythm, while you yourself show interest in harmony and improvisation lines. Is this attitude of holding interest something that all drummers have, or is it unique to you, Yoshijima-san?}

I simply think there are two patterns: people who originally played piano or guitar may be concerned about what's happening now with melody or harmony, or their ears automatically go there. People who started with percussion instruments, especially skin instruments, don't go there and listen at a completely different place.

\item [About the work's booklet, could you understand the content?]\mbox{}\\

There was nothing I couldn't understand. I could read it normally.\\

--- {\bf There were parts that weren't written in the usual notation, but you could still read it without any problems?}\\

It's just that you can read each part separately (so I could read it). The difficulty is when it flows simultaneously, but I can understand each individual thing that's written.\\

\item [About the booklet, were there any impressions or useful insights for you?]\mbox{}\\

For example, if there's one rail, and there's supposed to be another flowing in a different direction, but somehow they're on the same rail - that kind of sensation only comes out if you do music like that, and it became audible to me while performing. So I've already been able to apply it within myself, and I can bring it out in other situations, which is the cause of recent incidents.\\

--- {\bf It's amazing that you're using it in the field.}\\

When I started using it, nobody understands, after all. But depending on the person, they challenge on their side, so some bravely say ``keep it as it is.'' It's in a state where (beat) counting is not allowed, or rather, you can't count from the second beat onwards. Since I'm used to it, I can perform without a running start. I've gotten used to the state of dividing into 4 while deviating from the four beats that I was counting ``one, two, three, four'' until just before, so I can change it within myself. When asked to ``do an open solo,'' I do it ``intending to make them regret it.''\\

--- {\bf In an open solo, you can't tell where the downbeat is, right?}\\

I've also gone to that side, so I can't go back. You can't discard what you've learned, right? Once you get used to it, you want to try it in the field too.\\

\item [Have you experienced any similar music before?]\mbox{}\\

Similar examples to this are difficult. As I pointed out during the recording, if you write Bitches Brew in a score and try to create that unique atmosphere, you might need a score like this. But I've never seen the opposite pattern like this time. I hadn't yet had the experience of it being written in a score, and as a result of performing it, it became like Bitches Brew. However, by doing this kind of score this time, there might be something close to this in what Mr. Kikuchi was doing, or what Tipographica was doing.\\

--- {\bf I think knowing or not knowing that notation changes the analyzability of the music you listen to. If you write the notation like that, you might think ``maybe this music can be analyzed,'' but in the case of polyrhythm/cross-rhythm, sharing the bar line is a prerequisite, and in the case of wrap-around, sharing the beat is a prerequisite - it seems like there's only one or the other, and this time you're doing both.}\\

For me, there's surprisingly neither, and I read and perform melody first, so I only think about where to breathe. For instance, if there's music in 5 beats, even if the music's rhythm is in 5/4, if the breathing place is ``one, two, three, four, five, six'' at the end, I perceive that as one unit. So actually I might not be counting time signatures. Everything is in 1 beat.\\

--- {\bf That's like Indian musicians (in the sense of accumulating).}

So, I've performed Frank Zappa's Black Page before, but I don't understand what's difficult about it. You just have to read everything. The hi-hat is playing quarter notes. I understand (technically) that the difficult part is the many tuplets, but since you just have to read it, I've never thought it was difficult.\\

--- {\bf Music like Black Page is a virtuosic piece of contemporary music style. Compared to that difficulty, if you just look at the lead sheet of the piece you performed this time, anyone could perform it. But in reality, everyone said it was difficult, difficult when performing. What do you think is the difference in difficulty?}\\

It's that this time, it's not that each individual line is difficult, but the ensemble is difficult. For example, if three people have to do it, you hear completely different things, and within that, you have to maintain the pulse you need to put out. On top of that, if we compare with Black Page, Black Page has the quarter note decided. The performers ride on top of that to perform, so the difficulty of Black Page is that it's difficult for everyone to unison such odd tuplets together. So if you do it alone, it's not difficult.\\

--- {\bf If you feel like singing, you can sing Black Page by yourself.}\\

Because you just have to read what's written. But this time's piece, even if you read what's written, one line alone is meaningless. So it's an etude for ensemble practice, not for individual practice. However, what a drummer does to make it work as an ensemble is to hold the other parts together. So in the middle, how to make those misaligned beats into one within your body - this kind of solo ensemble is the first time.\\

--- {\bf The talk of holding together came up, but from an outside perspective, I thought Yoshijima-san was the one who understood best what the music was doing among the three. The other two were saying things like they didn't quite understand ``what's correct,'' but only Yoshijima-san could say ``that was correct now'' - a situation arose, and I think there's a difference in the role of the parts.}

\item [Is there any difference compared to what you usually do?]\mbox{}\\

Conversely, situations where I'm required to play normally are more difficult. When each person is only doing their own thing, even if I try to connect, it's impossible to hold it together.\\

--- {\bf Is it easier this time when everyone is trying to connect?}\\

Yes. That's ultimately only about whether it's musical or not. In that sense, this piece was more comfortable for me emotionally. No matter how high the individual abilities are, situations where people who play selfishly gather are more painful.\\

--- {\bf This time, watching the recording, the performance didn't rush (the tempo didn't gradually increase), but I thought maybe this music doesn't allow such things. Because everyone is trying to maintain the structure. Of course, there are higher levels, so I can't say that definitively based on this time alone, but what do you think?}\\

As we get used to the piece, we naturally want to challenge. I think rushing or dragging happens with that challenge image.\\

\item [Where do you think the difficulty of live performance lies?]\mbox{}\\

I think the most difficult thing is that you have to ``solo'' after grasping that. I think it would be good if someone could be the foundation when someone takes a solo, but this time the drums had a large proportion of that, but the difficulty changes depending on who puts out that pulse. For example, if the drums start soloing, who do you entrust the rhythm to? I would probably destroy it with considerable momentum, and by doing so, I think the difficulty increases further.\\

\item [How do you think about the development potential of this work and method?]\mbox{}\\

It's difficult, maybe it already existed. This time it's a special case of creating it from a score, but surprisingly when you listen to it, there was this ``this sounds like this,'' so maybe there won't be visible changes, but the fact that it can be controlled with a score is frightening. For example, things that were randomly entered in electronica like ``it's randomly entered but feels good'' - now it's controlled. I think there's tremendous progress in terms of reproducibility. The strength that this can be heard again. And it's not left to computers but performed by humans. This is frightening.

--- {\bf Is the point of the current discussion about reproducibility through scores?}\\

People like me who are eccentric are pleased. When playing randomly, you say ``I wasn't in good condition today,'' right? Like, we didn't match well. That disappears - I feel like stress is removed. Because it's visible that if you do this, you can produce the sound you want to produce - even if it's a minority, it really helps me. I don't have to fight anymore.\\

--- {\bf It's amazing that Yoshijima-san is developing or inheriting this rhythm in the field and using it. Is there anything about this entire project that hasn't come up yet?}\\

The performances I've done until now, even in my mind, are not very clear, and I wonder where my ambidexterity comes from and where it lands - this was one answer to that. But no one could systematize it, and even in my own head, it wasn't something that could become systematic. But what's interesting is that even though I was doing it before, when I was given the score and tried to do it, I couldn't. I got it to fit perfectly just before the recording, and I thought ``this isn't much different from the performance I always do.''\\

--- {\bf I think this work, rather than strongly asserting novelty, merely shows one way that occurs simultaneously and spontaneously as popular music?}\\

What's happening now is the same as what Berklee did with what Charlie Parker was performing, and I think Parker might think that if his performance hadn't been analyzed, jazz wouldn't have developed so much.\\

--- {\bf Historically, jazz is highly analyzable music, and it spread widely because Berklee Method fit into it, right?}\\

That's right. Whether that analysis is correct or not, various ways of thinking came out based on that, didn't they? In some cases, as if to deny that, Miles started doing modal. But then it was chaos, and nothing more came out, so there's a background of going to free jazz, but actually it came out, like this. The fact that it has reproducibility is enormously significant. But whether anyone can do it is a different story, including whether they want to do it.
\end{description}

\subsection{Interview with Participating Musician 2 (Yuji Kawamoto, December 18, 2017)}

\begin{description}
\item[Impressions after watching the video and listening to the audio of this work]\mbox{}\\

The impression was that it had become proper music more than I thought. That's due to playability issues and also the content written in the score, and I thought ``will this be okay?'' but it became more comfortable to listen to than I expected, so that was good, and at the same time it was mysterious or interesting.

--- {\bf In what way was it interesting?}\\

Such rhythm figures are written with music theory matters as the main focus rather than musical matters first, right? Because it's research with a purpose, that's fine, but when we actually performed and recorded it, it became a good balance, and I don't know if that's just because our ears got used to it, or the work was originally good, or the performers' ability - I don't know which, but I'm very satisfied with what was completed.\\

--- {\bf Behind the words ``it was more musically viable than I thought,'' I think there's the thought that ``it might not work,'' so what kind of non-viable state were you thinking of?}\\

I was thinking that. I wondered if it would be okay. I understood the concept, but I thought it might be impossible to perform. That was the biggest thing. Next, because there are many research constraints, the arrangement of notes written there might not sound like melody, and if it doesn't sound like melody, it won't become music. But it sounds. Especially Etude \#1 - I thought ``will this be okay?'' but I think the performers' ability was significant.\\

--- {\bf What do you think about the criticism that might be received that ``the performers are just matching it''?}\\

This is the dilemma of such research - when the criticism ``aren't the performers just matching it?'' comes out, from a musical perspective, ``what's wrong with matching?'' Even if there's a physical problem of whether such an attempt is possible or impossible as a research theme, as a performer's final responsibility, we have to make something happen, so ``we couldn't do it'' or our pride of course, and the desire to contribute to research is also there, so from that standpoint difference, the sense of the musical field is ``what's wrong with matching?'' I think that contradictory position always emerges in doing such things. I think it's about how to twist that down as a work.\\

\item [About the work's booklet, could you understand the content?]\mbox{}\\

Of course I could understand. There's more than enough explanation.\\

\item [About the booklet, were there any impressions or useful insights for you?]\mbox{}\\

This won't directly answer the question, but from musicians' musical information perspective, I think there are parts that are slightly over-explained. It could be simpler. However, I thought the innovation of dividing the drum score into two staves was very good. I made my own programmed sound source for practicing at home, and such things were very good for understanding musical layers. It was very good that it showed that this plus this becomes this. As a score to come out in the field, it was something I'd never seen before, but I think it's sufficient as a project instruction manual. Also, the talk about cognitive dissonance was interesting. Very Yamato-like.\\

\item [Have you experienced any similar music before?]\mbox{}\\

I think there are many similar examples, and jazz in general is like that.\\

--- {\bf Generally, when poly is happening, the bar lines are shared, and phrase wrap-around has shared beats. If the characteristic of this work is that neither bar lines nor beats are shared, and what's shared is the pulse at the least common multiple with large blocks, then are there similar examples?}

That's difficult. There might be in classical works, but I can't say this specifically, but in jazz, such things can probably happen during improvisation, and I think Evans was probably doing it, and Tony Williams in Miles' band was probably doing it.\\

--- {\bf This time we used variable meter notation, but usually we don't analyze like that. Because the analysis method is done in the form of Western music or Berklee notation, we didn't understand, but if you analyze Tony Williams etc. in variable meter notation, you can imagine you'd find he was doing similar things.}

That's probably true. As Yamato-san says, not limiting to Berklee, using staff notation and bar lines itself is a Western system, a Western European folk music system, and for example, Indian music uses different notation, so when you look at it with such different vocabulary, it's not strange at all that there are different analysis methods or analysis methods most suitable for music with such structure.\\

--- {\bf When African poly is transcribed in Western musical notation, Western music misses African poly - this is written in Kenichi Tsukada's ``The Truth about African Music,'' and they're missing it because they don't have the perspective of what kind of bitmap or structure to view it as.}\\

In that case, from the perspective of doing jazz, the rhythm of ``swing'' is impossible to notate on staff notation, so I don't trust scores very much to begin with. Also, there are five types of beats in jazz, and great performers like Parker firmly control which position they're performing at for one pulse. Laid back is also laid back, but there are two types of laid back, and they control that. Why jazz became like that is because I think the African people's ethnic subtle off-score concept was inherited directly into jazz, and surprisingly, that's passed down orally, not in scores - that's the fact. I think that's something unique to Black people when considering jazz's ethnicity. Japanese people aren't that interested in continuous rhythm.\\

--- {\bf Don't you think there's something common with Robert Glasper's rhythm sense?}

I think that also came out from past context, and jazz musicians have high affinity because the roots are the same, so I think it came out from such context. I read Gunther Schuller's book ``Early Jazz,'' and Gunther Schuller writes that the subtle deviation of fine pulses in African ethnic rhythms is not accidental deviation, but jazz giants like Charlie Parker and Dizzy Gillespie were doing it while recognizing it.\\

\item [Is there any difference compared to what you usually do?]\mbox{}\\

I think this time can be said to be ``vocabulary research.'' Vocabulary is very important whether composing or playing, but beyond that is ``what do you want to say?'' - when making music, we do all of that comprehensively, but I think it's grammar research with focus on one point within that, and I think it's something that really should be done, but usually we can't spend time focusing on that one thing to that extent. That's why I understood the concept but didn't know if I could perform it. That's why I thought it was very valuable.\\

\item [Where do you think the difficulty of live performance lies?]\mbox{}\\

Even if you learn new words, to be able to use those words freely, you need to physically internalize them once. Isn't that where it is? In other words, ``it takes time.''\\

--- {\bf When physically internalizing, you seemed to struggle with 5 beats in 4/4?}\\

Because I don't usually experience it. What I understood this time is that I normally do 3 beats in 4 or 2 beats in 3, but why I can do it is because I usually do it. I've never thought about 3 beats in 4 with this kind of resolution this time. In other words, I could play it because I repeatedly listened normally, memorized, and played, so if I usually listened to and memorized songs in 5 beats in 4, I could probably do it.

--- {\bf In that sense, were the simulated demo sound sources useful?}\\

Of course. In the sense of physically installing, I can't say it was useful because I made practice karaoke myself, but for the demo sound source, the bigger issue is whether it sounds like music, so the surprise of ``so it becomes like this'' was bigger. After that, it was a feeling of processing it in the field.\\

\item [How do you think about the development potential of this work and method?]\mbox{}\\

When I saw this work, the first thing I thought of was Bartók's string quartet. This becomes more interesting if you apply it to counterpoint. Rather, I think Bartók and Stravinsky are doing it. But isn't it applicable anywhere? I thought I'd do such things in my own songs. So in that sense, I thought about analyzing it. If you only follow the vertical lines, you don't really understand what's happening, but when 5 phrases or 4 phrases coexist, the chords that happen incidentally become the music's coloring - you can make music with that way of thinking, right? If you decide on tonality or decide on scale. Then I thought it would be useful for counterpoint-like things, and it's good to do it during improvisation, and that's like what Yoshijima-san is doing in the field, but if you suddenly do it, everyone would be surprised.\\

--- {\bf It's similar to Miles' reference to Tony Williams.}\\

Even if it's a word thought out logically, if it doesn't convey to people, it's meaningless - or rather, if you suddenly start speaking words that don't make sense, there's the aspect that people around don't understand, so I think that's really difficult, but still, such things should be done, and not in improvisation but properly done within musical pieces.\\

--- {\bf Firmly holding it as structure?}\\

Yes.\\

--- {\bf If asked about this research, it's not about novelty or anything, but I think ``the value is in describing such structures, making them viable as music, and performing them live''?}\\

I think this is like Matsuo Bashō. In other words, it's ``research on new vocabulary.'' It's ``research on new expression methods.'' It's not a matter of novelty or such problems, it may be new, but.\\

--- {\bf Popular music occurs contemporaneously, so asserting novelty doesn't make much sense.}

Going as far as a new format might be too much, but for example, if you divide 5 beats into 4, you can express fuzzy fluctuations that couldn't be expressed with sixteenth notes before, right? Based on that, there's the possibility of creating music with different expressions that's more seamless. Research on that vocabulary - it's like not just blue or red but purple, or yellow mixed into purple - such expressive range expands, right? If you can properly hold that in your physical being and use it freely, that's an enormously big thing. So it's like Matsuo Bashō - fitting your words and sensations into 5-7-5, though Bashō's vector is opposite, if you research like that, Bashō doesn't need to speak sentences. Such expression became possible. So it's the same - if you do more refined research, the expressive range will expand more. I think everyone thinks so, but when active jazz musicians do it, it comes down to how much value they find in terms of time. When you remove the cost-performance issue, there's nothing more interesting than this.
\end{description}

\subsection{Interview with Participating Musician 3 (Masaaki Kaneko, December 19, 2017)}

\begin{description}
\item[Impressions after watching the video and listening to the audio of this work]\mbox{}\\

Nothing in particular, but normally like ``we did it,'' ``we made it this far,'' ``amazing.'' The feeling of ``amazing'' from the side that performed. I often create from zero to one, so I understand very painfully how hard it is to do this far. In that respect, rather than how the research is, I think it's wonderful that ``such a result was produced.'' Moreover, considering the detailed scheduling and path to get there, it's heavy. You can think anything or express anything, but getting there is tremendously difficult. As for the evaluation of live performance as a research experiment, it's not for me to say, but with theory-oriented experiments like this, it's almost impossible to tell if it's useful or not, and there are many such parts, so there's no way to evaluate. Conversely, not at that level, but I'm evaluating that it was made, and there are almost no such people. Jazz musicians too make their own CDs, and when they think something, they can take action, but they don't actually do it. The fact that it was actually done is amazing.\\

\item [About the work's booklet, could you understand the content?]\mbox{}\\

If I read it with great determination, what they want to do would probably be conveyed, but just reading it superficially doesn't convey it. Since I haven't read it carefully, it hasn't entered me yet. I can't judge yet what it would be like when I read it carefully.\\

\item [About the booklet, were there any impressions or useful insights for you?]\mbox{}\\

Regarding rhythm, there certainly is much to think about, and human rhythm cognition was a point of interest.\\

\item [Have you experienced any similar music before?]\mbox{}\\

There are simpler ones. But not with 5 beats and 4 beats. I haven't crossed but made it poly, but in my own musical experience, I've practiced and used in improvisation things like dividing into several, but it doesn't go to poly. There's 3 or 4, but I've never done it that far. But there's nothing like mixing poly and cross.\\

\item [Is there any difference compared to what you usually do?]\mbox{}\\

I don't fully understand what the purpose of cross and poly this time is, but what I usually do is, to put it flatly, a business approach, and that's the difference.\\

--- {\bf In that sense, this relationship is also a business relationship of requesting and completing performance of a business piece. On top of that, I thought I could see the pride as a performer of ``if we can get this far, it's complete''?}\\

That's also there. It's a part I'm personally imposing on myself, probably a part I can't personally forgive. But if the director or producer says ``OK,'' there was also that aspect of having to stop there. That inevitably exists. But when involved, I think there's such a growth instinct.\\

\item [Where do you think the difficulty of live performance lies?]\mbox{}\\

It's because the rhythm is compound, right? Simply that. To put it simply, I thought multitasking is difficult. I've heard that when humans do two things, they can't do 50% each, but about 30%, which is inefficient. I think it's that kind of difficulty. However, I feel that when humans process things like this time, they're not doing multitasking. Isn't it close to drummers not processing with multitasking? Like grasping it vertically or something.\\

--- {\bf Grasping it as one wave like intricate knitting?}\\

Yes. Probably, because we don't have that kind of learning of grasping method until now, we have trouble performing things like this time. I think if you learn that, you gradually get used to it and people who can do it can do it. How are humans processing things like this time? I'm very interested in that perspective. From that perspective, the act of performing is very intellectual or artistic work.\\

\item [How do you think about the development potential of this work and method?]\mbox{}\\

Couldn't people in the world think about how to use such music? At the time of recording, I also said this music would be good ``to play in bookstores,'' but I think ``music assimilates to places,'' so it would be good if it could be created creatively while assimilating to places.\\

--- {\bf Are you seeing ahead of the development of this music that there must be places suitable for this music?}\\

That's right. Not limited to this music, but art or anything, if not sought by people, it can't exist and just disappears, so there must be a good place somewhere for what was born. If music is in the flow of BGM, I think it assimilates to places, so I think it's important for artists to think about where that music is needed.
\end{description}

\subsection{Interview with Musicians from Other Genres (Shinya Ando, Jun Haraguchi, December 13, 2017)}

\begin{description}
\item[Impressions after watching the video and listening to the audio of this work]\mbox{}\\

{\bf Ando} In Etude \#2, seeing the bassist counting 5 with his mouth, I understood ``oh, it's 5'' - that was my first entry point. From there I understood that it can be taken as either. As a piece, it was very good. Because the rhythm division is unique and strange, there was a unique groove feeling and it was very interesting. It was cool. Melodic instrument players, including myself, do the difficult thing of forcefully putting in odd meter and keeping it while playing solos, but it's still difficult. That's also toward the future, right? When an era comes where there are lots of standard songs in 5/4 first, and people are used to 5, it might gradually become possible to put in different rhythms quickly. In parts like solos.\\

{\bf Haraguchi} Etude \#2 had a feeling that drums could be taken in 4 or taken in 5, and if you change your way of listening, it's easy to hear either way, and as musicians, I was listening wondering if that makes it easier to perform. The part where bass and piano switch rhythms in Etude \#2 is good as a piece. The feeling of the music's axis changing sharply is interesting. The fact that the bass plays four quarter notes there was good. Rather than doing various things melodically, such a quarter note line here emphasizes the change, so the feeling of change comes strongly and it was interesting. If you just listen to the sound without hearing any scores or explanations, you don't understand what's being done just from that - you can't analyze/parse it, but as a piece it was good.\\

\item [About the work's booklet, could you understand the content?]\mbox{}\\

{\bf Haraguchi} I understood the content, but listened to grasp the image and was convinced. I was reading, looking at the booklet, then while listening and looking at the score, I could be convinced like ``oh I see, I see.''\\

{\bf Ando} Since I talked with the work's creator (in advance), I have that preconception, so I don't know if I could understand when seeing this without knowing anything at all.\\

\item [About the booklet, were there any impressions or useful insights for you?]\mbox{}\\

{\bf Haraguchi} With instruments like santoor or tabla where both hands can be used, it's possible, but there are various examples of two performers performing in different rhythms, not like what drums are doing here of performing different beats simultaneously with left and right hands. However, there was no past approach of one performer doing different beats with both hands. Ando-san said in summer that when he met Yamato-san, things he was vaguely thinking about took shape, and he's actually doing it recently.\\

{\bf Ando} I'm thinking about how to do it, but since it's still in a basic state, I'm only doing basic things like normally right hand is 7 and left hand is 4. When the base is 10 beats, I divide in the middle and put 7 and 4 against those 5 beats.\\

{\bf Haraguchi} At that time, with tabla keeping 5, one hand doing a 7 phrase and one hand doing 4 is done melodically, so the approach of doing it where it sounds like an interesting melody is what we've been doing recently.\\

--- {\bf When dividing that 10 by 5, you're doing polyrhythm in a state where the bar as a certain unit is shared, which is a similar way of thinking to Afro-poly, right?}

{\bf Ando} It's a way of thinking like triplets in 2 beats.

\item [Have you experienced any similar music before?]\mbox{}\\

{\bf Haraguchi} This is an example of Zakir Hussain and Mickey Hart performing in '76, but it's an example of taking it in a very large span. It doesn't correspond to this example, but reading the explanation of the piece, there's an example of making one melody within 102.5 beats. However, just listening to this piece, you can't take it as a 102.5 melody, but according to the explanation, it's like that. This is Indian classical music, done by tabla player Zakir Hussain. He's performing with various percussion instruments.\\

{\bf Ando} What Zakir Hussain is doing is arranging and changing pieces, but the way of thinking is only Indian music. Even in Shakti, he's changing various things but only doing Indian classical.\\

--- {\bf Even listening to this, you don't understand what unit the length of one beat is, right?}

{\bf Ando} You don't understand, right.\\

\item [Where do you think the difficulty of live performance lies?]\mbox{}\\

{\bf Ando} The difficulty is, if you present something like a table in the booklet, there's no difficulty, but continuing it to the speed you want to play takes a long time, right? It's not at the level of practicing one phrase thinking ``this is a bit difficult.'' It takes so long that you wonder if you'll reach it, if you try to do it alone. Doing it as a group is a bit easier, and if you remember how to fit it, if it just pops into you. In return, if you have in your hand the tiles you're putting out, like in 5/4 it becomes like this, or 7 is like this, 3 is like this, 6 and 12 are like this, you can fit them in, but for one person to start doing separate things like this drummer takes time.\\

--- {\bf When doing such performance and actually doing ensemble, do things happen like being pulled in?}\\

{\bf Ando} You get pulled in.\\

{\bf Haraguchi} You're about to get pulled in.\\

--- {\bf How do you maintain not being pulled in?}\\

{\bf Ando} The feeling of becoming a metronome yourself.\\

{\bf Haraguchi} I'm singing rhythm within myself, but the basic rhythm in my heart. While singing, I'm listening to the other side. Looking at the head, of that melody's point.\\

{\bf Ando} Looking at the points, and since they're different, I have the image of going around like this.\\

--- {\bf Regarding the point of how much you're listening to the other person in ensemble, how is it for you two? I think it's hard for a state of listening to the entire ensemble to occur, focusing on necessary parts at times?}\\

{\bf Haraguchi} As for what drums are listening to, there's a similar role. The side singing melody listens to sounds that are landmarks.\\

{\bf Ando} In polyrhythm, for example, in 10/4, the head and the 6th beat and the back of the 3rd beat, I'm also concerned about the 3rd beat. I'm concerned about those, but not much else.\\

--- {\bf Is it a feeling of doing ensemble with structural understanding of rhythm mutually?}\\

{\bf Ando} That's right. This method itself is rhythm, right? Because it's something that can stand with just rhythm, we're putting melody on top of that, so rhythm awareness is very high, especially in India. I think there are many musicians who aren't like that, but listening to Indian people's performance, rhythm awareness is high, right? There's also a way of putting 3 into 1 beat, like shifting by 5.\\

{\bf Haraguchi} Like where you are, in the blank of the beat. At 1/3 of the beat or 1/5 of the beat. I'm conscious of that, or rather, I try to be.\\

{\bf Ando} When doing polyrhythmic things, I'm listening to the other person very much. When doing polyrhythmic trick-like things more than usual, because the points have to match, I'm listening considerably. The side doing the base beat gets pulled in, so while listening to the other person, they do it like a metronome, but the side approaching is listening more. Listening while fitting phrases.\\

--- {\bf Listening to the current talk, rather than harmonic rhythm, it's really with ``rhythm'' that you're in a state of being able to grasp where your position is?}\\

{\bf Haraguchi} That's right.\\

{\bf Ando} Of course, we use that thing called harmonic rhythm, but before that, rhythm as a thing is (told by my teacher). Ravi Shankar's teacher's son is there, and that person has mouth in some meter, hands in some meter, mouth and both hands, both feet are separate in about 5 or 6 things.\\

{\bf Haraguchi} (I) can do up to 3. I can do different things with mouth and both hands. But when it comes to both feet (it's difficult). Like the sensation of shifting in - a characteristic of Indian classical music is that the melody line of a unit straddles and shifts across this bundle...\\

{\bf Ando} There's clearly going with 7 in a 10-beat piece.\\

--- {\bf Sharing 7 beats while sharing the beat in 10 beats and it shifts.}\\

{\bf Ando} There's also shifting like tihai.\\

--- {\bf At that time, the vertical line beats are aligned.}\\

{\bf Ando} Yes. There's also that.\\

\item [How do you think about the development potential of this work and method?]\mbox{}\\

{\bf Haraguchi} For me concretely, I thought there are elements of this that can be incorporated into my performance.\\

{\bf Ando} Of course there are, and I want an era to come when music with such multilayered rhythm is playing when you walk around town normally.\\

{\bf Haraguchi} Because very simple music is rampant in the world, it's boring or rather, only things that anyone can immediately understand at first glance are playing, but you can't understand analysis properly without facing this carefully, right? You don't understand what's happening, but when you listen to it as music, it's very beautiful - I want such music to spread more.\\

{\bf Ando} I forget why it's 4, but in other words, without such art music, a state where the masses who listen are strong, even in economic terms, in the past there were patrons who employed musicians, and I think music grew from that. That's sponsors putting it out, or now it's an era of selling tickets and from there - isn't it a structure where the listening side is strong? I've read somewhere that when the masses choose, music and such things become increasingly low-intelligence, and when touring around, that's actually the case. Rather than doing interesting things, I don't do them. There are many people with the feeling that unless it's something satisfying within their small skill set, they won't recognize it. If it gradually becomes like that, the world becomes boring, right? So I imagine a future where such music spreads more. I want it to be playing at friends' shops.
\end{description}

% References
\bibliographystyle{chicago}
\begin{thebibliography}{99}
\bibitem[Murai 2017]{murai:01} Yasushi Murai 2017 "Jazz History That Changes Your Way of Listening" Shinko Music Entertainment Co., Ltd.
\bibitem[Harari 2016]{harari:01} Yuval Noah Harari (translated by Hiroyuki Shibata) 2016 "Sapiens: A Brief History of Humankind (Complete Edition: The Structure of Civilization and Human Happiness)" (Kindle Locations 3000-3010) Kawade Shobo Shinsha
\bibitem[Van der Merwe 1999]{merwe:01} Peter van der Merwe (translated by Toyo Nakamura) 1999 "The Basic Theory of Popular Music" Music Magazine Co., Ltd.
\bibitem[Kikuchi \& Otani 2004]{kikuchi-ootani:01} Naruyoshi Kikuchi, Yoshio Otani 2004 "The School That Taught Melancholy and Sensuality" Kawade Shobo Shinsha
\bibitem[Asaka 1977]{asaka:01} Jun Asaka 1977 "New Music Dictionary" Ongaku no Tomo Sha
\bibitem[Yoda \& Ono 2016]{yoda-ono:01} Sho Yoda, Takashi Ono 2016 "Music Theory Analysis of Polyrhythm Typology" Shinshu University Faculty of Education Research Journal No. 9 pp.169-188
\bibitem[Bosseur 2008]{bosseur:01} Jean-Yves Bosseur (translated by Utako Kurihara) 2008 "88 Keywords to Decode Contemporary Music" Ongaku no Tomo Sha
\bibitem[Tsukada 2016]{tsukada:01} Kenichi Tsukada 2016 "The Truth about African Music" Ongaku no Tomo Sha
\bibitem[Deva 1994]{deva:01} B.C. Deva (translated by Hiroshi Nakagawa) 1994 "Introduction to Indian Music" Toho Shuppan
\bibitem[Sachs 1979]{sachs:01} Curt Sachs (translated by Shigeo Kishibe) 1979 "Rhythm and Tempo" Ongaku no Tomo Sha
\bibitem[Mahdi 1998]{mahdi:01} Salah al-Mahdi (translated by Yoshiko Matsuda) 1998 "Arab Music" PASTORALE Publishing
\bibitem[Clave 2008]{clave:01} Rumba Clave. 2008 "Rumba Clave: An Illustrated Analysis"\\
\url{http://rumbaclave.blogspot.jp/2007/08/rumba-clave-illustrated-analysis.html}: Accessed January 3, 2018 (Wed)
\bibitem[Obrist 2015]{obrist:01} Hans Ulrich Obrist (translated by Naoko Shinogi, Fumiko Uchiyama, Hisashi Nishihara) 2015 "Music: Composers Who Created 'Contemporary Music'" Film Art Inc.
\bibitem[Strickland 1998]{strickland:01} E. Strickland (translated by Toshie Kakinuma, Sakae Yoneda) 1998 "American New Music" Keiso Shobo
\bibitem[Kikuchi 2003]{kikuchi:01} Naruyoshi Kikuchi 2003 "Complete Track Commentary by Naruyoshi Kikuchi Himself on New Album 'Structure et Force'!" \url{https://web.archive.org/web/20031002092226/http://www.bls-act.co.jp/artists/dcprg.html}: Accessed January 8, 2018 (Mon)
\bibitem[Yubi 1996]{yubi:01} Kuniko Yubi 1996 "All About Popular Rhythms" Keiso Shobo
\end{thebibliography}

\begin{referenceworks}{99}
\bibitem[Yamato 2017a]{yamato:01} Hiroshi Yamato 2017a "Music for multilayered temporal structures" \url{https://goo.gl/o1pTEk}
\bibitem[Yamato 2017b]{yamato:02} Hiroshi Yamato 2017b "4by5" \url{https://github.com/dropcontrol/4by5}
\bibitem[Yamato 2017c]{yamato:03} Hiroshi Yamato 2017c "3by4by5" \url{https://github.com/dropcontrol/3by4by5}
\bibitem[Yamato 2017d]{yamato:04} Hiroshi Yamato 2017d Documentary Video "Music for multilayered temporal structures" \url{https://youtu.be/3AuNz4GiLsI}
\bibitem[Yamato 2017e]{yamato:06} Hiroshi Yamato 2017e "Test6 (Piano Etude for D Dorian)" \url{https://soundcloud.com/dropcontrol/test6piano-ensemble-for-d-dorian}.
\bibitem[Yamato 2017f]{yamato:07} Hiroshi Yamato 2017f "Test5 (Percussion Ensemble)" \url{https://soundcloud.com/dropcontrol/test5-percussion-ensemble}.
\bibitem[Yamato 2016]{yamato:05} Hiroshi Yamato 2016 {\it New York 2997} "PCM15 BEAT MUSICS/Penguin Music Graduate School 2015 Black Music Class Graduation Production Collection" Sony Music Artists Inc.
\bibitem[Davis 1959]{miles:01} Davis, Miles 1959 {\it So What} "Kind of Blue" Columbia
\bibitem[Davis 1970]{miles:02} Davis, Miles 1970 {\it Bitches Brew} Columbia
\bibitem[Davis 1956]{miles:03} Davis, Miles 1956 {\it Bye Bye Blackbird} "Round About Midnight" Columbia
\bibitem[Reich 2013]{reich:01} Reich, Steve 2013 {\it Piano Phase} "Steve Reich Early Works" Warner Music Japan
\bibitem[Reich 2016]{reich:02} Reich, Steve 2016 {\it Violin Phase} "The Ecm Recordings" Ecm Records
\bibitem[Reich 1995]{reich:03} Reich, Steve 1995 "Drumming" Nonesuch
\bibitem[Date Course Pentagon Royal Garden 2001]{works-kikuchi:01} Date Course Pentagon Royal Garden 2001 {\it PLAYMATE AT HANOI} "REPORT FROM IRON MOUNTAIN" P-Vine Records
\bibitem[Date Course Pentagon Royal Garden 2003]{works-kikuchi:02} Date Course Pentagon Royal Garden 2003 {\it structure I la structure de la magie monderne / Structure 1 (Structure of Modern Magic)} "Structure et Force" P-Vine Records
\bibitem[Glasper 2012]{Glasper:02} Glasper, Robert 2012 "Black Radio" Blue Note Records
\end{referenceworks}

% To fit page breaks and figures into chapters, insert page breaks as follows.
\afterpage{\clearpage}
\newpage

\section*{Acknowledgments}

In conducting this research, I have been greatly indebted to the following people. First, I deeply appreciate Professor Masahiro Miwa, Professor Shinjiro Maeda, and Professor Masahiro Kobayashi for their various advice and guidance over two years. I especially thank Professor Masahiro Miwa for giving me the opportunity to deeply consider "how to face music in my life hereafter." I also thank the jazz musicians who participated in the master's work: Tomohito Yoshijima, Yuji Kawamoto, Masaaki Kaneko, engineer Naoki Hayamizu, Takayoshi Oishi who filmed the documentary video, and Masaki Okubo who performed mixing and mastering. I especially received various advice and assistance from Takayoshi Oishi during the research process.\\

I thank Professor Yasuhiko Ando of ATP, the members of NxPC, especially Professor Mami Hirabayashi and Ayumu Nagamatsu. Thanks to you, I was able to realize several activities I couldn't do before. I also thank Associate Professor Shigeru Matsui, who provided advice from various angles as a critic. I still haven't answered the question "What is a masterpiece for you, Yamato-san?" that Associate Professor Matsui asked me by email last summer, but I have an answer prepared, so I would be happy if we could discuss it again at the next opportunity.\\

I thank Masayuki Yamabe, Saki Koga, Daiji Hirata, and Hiroshi Yamada, who pushed my back when entering IAMAS. Especially Yamabe, Koga, and Hirata have been my friends and mentors for over 10 years. I also send the same thanks to my colleagues at Alliance Port at that time and everyone who was involved there.\\

I dedicate this thesis to both Ginji Ogawa and Naruyoshi Kikuchi. I was an unworthy disciple of Ginji-san, but I learned various things about playing guitar. I also thank "Penguin Music University," Mr. Kikuchi's private school, and its members. Especially to Kohei Kanayama and Seiichiro Honjo. What I learned with them under Mr. Kikuchi was the beginning of all this research.\\

I dedicate my gratitude to my father, mother, and brother who always support me. I truly think the words my mother once said, "I'm so glad you're doing music. If you weren't, I think you would have died long ago," aptly represent my current self.\\

Finally, as this thesis shows, music is always an art that exists in relationship with someone and is "something that cannot be done alone." Therefore, I note at the end that there are many people I couldn't write here, and let this serve as acknowledgment to them.\\

Thank you all very much.

\end{document}
